\documentclass[11pt]{article}
\usepackage[a4paper,margin=1in]{geometry}
\usepackage{amsmath,amssymb}
\usepackage[T1]{fontenc}
\usepackage{lmodern}
\usepackage{microtype}
\usepackage{xcolor}
% Footnotes inside \parbox (used below) are otherwise trapped; savenotes emits them
% at the end of the wrapper so they appear on the page foot as usual.
\usepackage{footnote}

\setlength{\parindent}{0pt}
\setlength{\parskip}{0.5em}

\newcommand{\promptbox}[1]{%
  \begin{savenotes}%
  \begingroup\setlength{\fboxsep}{10pt}\noindent
  \colorbox{blue!8}{\parbox{\dimexpr\linewidth-2\fboxsep\relax}{\textbf{Prompt. }#1}}
  \endgroup
  \end{savenotes}%
}
\newcommand{\resultbox}[1]{%
  \begin{savenotes}%
  \begingroup\setlength{\fboxsep}{10pt}\noindent
  \colorbox{purple!8}{\parbox{\dimexpr\linewidth-2\fboxsep\relax}{\textbf{Result. }#1}}
  \endgroup
  \end{savenotes}%
}
\newcommand{\examplebox}[1]{%
  \begin{savenotes}%
  \begingroup\setlength{\fboxsep}{10pt}\noindent
  \colorbox{green!8}{\parbox{\dimexpr\linewidth-2\fboxsep\relax}{\textbf{Example. }#1}}
  \endgroup
  \end{savenotes}%
}

\title{\textbf{Book 1 --- Lecture 2 --- Exercise 2}\\[2mm]
\large Second Derivatives of Four Introductory Functions}
\author{}
\date{}

\begin{document}
\maketitle

\promptbox{The derivative of a derivative is called the second derivative and is
written \((d^2 f(t))/d^2\).\footnote{Operationally, the second derivative is ``the
derivative of the derivative'': one first computes \(f'(t)\), then differentiates
that new function of \(t\) again. Standard Leibniz notation writes the variable in
the denominator twice, e.g.\ \(\dfrac{d^2 f}{dt^2}\) for a function of \(t\), and
\(\dfrac{d^2 g}{dx^2}\), \(\dfrac{d^2\theta}{d\alpha^2}\), \(\dfrac{d^2 x}{dt^2}\)
for the other independent variables below. The \(d^2\) in the numerator indicates
that differentiation has been applied twice.} Take the second derivative of each of
the functions listed above:}
\[
\begin{aligned}
f(t) &= t^4 + 3t^3 - 12t^2 + t - 6, \\
g(x) &= \sin(x) - \cos(x), \\
\theta(\alpha) &= e^\alpha + \alpha\ln(\alpha), \\
x(t) &= \sin^2(t) - \cos(t).
\end{aligned}
\]

\section*{\(f(t)\)}
From Exercise 1, \(f'(t)=4t^3+9t^2-24t+1\).\footnote{Exercise 1 computed first
derivatives of the same four functions.} Differentiate again:\footnote{For a
polynomial, each differentiation lowers the highest power by one; here \(f\) was
degree \(4\) in \(t\), so \(f'\) is degree \(3\) and \(f''\) is degree \(2\).}
\[
f''(t)=\frac{d^2f}{dt^2}=12t^2+18t-24.
\]

\section*{\(g(x)\)}
We had \(g'(x)=\cos(x)+\sin(x)\).\footnote{As in Exercise 1, angles are in radians.}
Thus
\[
g''(x)=\frac{d^2g}{dx^2}=-\sin(x)+\cos(x).
\]\footnote{Apply \((\cos x)'=-\sin x\) and \((\sin x)'=\cos x\) to \(g'(x)\).}

\section*{\(\theta(\alpha)\) (\(\alpha>0\))}
We had \(\theta'(\alpha)=e^\alpha+\ln(\alpha)+1\). Differentiating once more term by
term,\footnote{For \(\alpha>0\), \(\frac{d}{d\alpha}\ln(\alpha)=1/\alpha\). The summand
\(+1\) is constant and has derivative \(0\). Also \((e^\alpha)'=e^\alpha\).}
\[
\theta''(\alpha)=\frac{d^2\theta}{d\alpha^2}=e^\alpha+\frac{1}{\alpha}.
\]

\section*{\(x(t)\)}
One convenient form of the first derivative is \(x'(t)=\sin(t)\bigl(2\cos(t)+1\bigr)\).
Differentiate using the product rule:\footnote{Write \(x'(t)=u(t)v(t)\) with
\(u=\sin t\) and \(v=2\cos t+1\). Then \(x''=u'v+uv'\).}
\[
\begin{aligned}
x''(t)=\frac{d^2x}{dt^2}
&=\cos(t)\bigl(2\cos(t)+1\bigr)+\sin(t)\bigl(-2\sin(t)\bigr)\\
&=2\cos^2(t)+\cos(t)-2\sin^2(t)\\
&=2\cos(2t)+\cos(t),
\end{aligned}
\]
using \(\cos^2 t-\sin^2 t=\cos(2t)\).\footnote{This is the cosine double-angle identity;
it rewrites the combination \(2\cos^2 t-2\sin^2 t\) as \(2\cos(2t)\).}

\resultbox{In summary,
\[
\begin{aligned}
\frac{d^2f}{dt^2} &= 12t^2+18t-24,\\[2pt]
\frac{d^2g}{dx^2} &= \cos(x)-\sin(x),\\[2pt]
\frac{d^2\theta}{d\alpha^2} &= e^\alpha+\frac{1}{\alpha}\quad(\alpha>0),\\[2pt]
\frac{d^2x}{dt^2} &= 2\cos(2t)+\cos(t).
\end{aligned}
\]}

\examplebox{At \(t=0\): \(f''(0)=-24\), \(x''(0)=2+1=3\). At \(\alpha=1\):
\(\theta''(1)=e+1\).}

\end{document}
